\chapter{INTRODUCCIÓN}

% ------------------------------------------------------------
\section{Contexto, planteamiento del problema y justificación}
% ------------------------------------------------------------

En los últimos años, el Internet de las Cosas (IoT, \textit{Internet of Things}) ha cambiado
la manera en que se usan las redes de comunicaciones. A diferencia del tráfico tradicional,
donde unos pocos usuarios generan grandes volúmenes de datos, las aplicaciones IoT se
caracterizan por millones de dispositivos que envían paquetes pequeños de forma esporádica
y sin coordinación previa~\cite{ref1, ref2}. Se estima que para 2029 habrá más de 29 mil
millones de dispositivos IoT conectados, impulsados por casos de uso como ciudades
inteligentes, agricultura de precisión, monitoreo ambiental y transporte
conectado~\cite{ref3}. Este crecimiento plantea un reto nuevo para las redes celulares,
porque los mecanismos de acceso pensados para comunicaciones humanas no fueron diseñados
para soportar densidades tan altas de dispositivos.

Para atender este escenario, el proyecto 3GPP (\textit{3rd Generation Partnership Project})
introdujo en su Release 13 el Internet de las Cosas de Banda Estrecha (NB-IoT,
\textit{Narrowband Internet of Things}), una interfaz de
acceso radio optimizada para comunicaciones tipo máquina (MTC, \textit{Machine-Type
Communication}), que opera en espectro licenciado con un ancho de banda de
180\,kHz~\cite{ref3, ref7}. Aunque NB-IoT permite soportar una gran cantidad de dispositivos
de baja complejidad, su canal físico de acceso aleatorio (NPRACH, \textit{Narrowband Physical
Random Access Channel}) dispone de solo 48 preámbulos ortogonales por oportunidad de
acceso~\cite{ref7}, lo que se vuelve un cuello de botella cuando cientos o miles de
dispositivos intentan conectarse al mismo tiempo. En esas condiciones, la probabilidad de
que dos o más dispositivos escojan el mismo preámbulo crece rápidamente, se producen
colisiones y aparecen retransmisiones que aumentan el retardo de acceso y reducen la
capacidad efectiva del sistema~\cite{ref4, ref5}. La literatura coincide en que los
esquemas de acceso ortogonal convencional no son suficientes para soportar las densidades
de dispositivos que se proyectan en escenarios IoT masivos~\cite{ref4, ref5}, y por eso el
problema del acceso aleatorio en NB-IoT sigue siendo un tema activo de investigación.

A esta limitación de capacidad se suma una limitación de cobertura. La infraestructura
celular terrestre llega principalmente a zonas urbanas, dejando sin conectividad confiable
a regiones rurales, marítimas y remotas, que son justamente donde muchas aplicaciones IoT
operan de forma permanente~\cite{ref3}. Aplicaciones como el monitoreo de cultivos, el
rastreo de flotas en zonas marítimas, la vigilancia ambiental en áreas protegidas o el
soporte a comunidades rurales dependen de que la red sea capaz de atender muchos
dispositivos con bajo consumo y cobertura global, algo que solo se logra combinando
tecnologías de baja potencia como NB-IoT con infraestructura satelital~\cite{ref3, ref7}.
Esta brecha ha motivado el uso de redes no terrestres (NTN, \textit{Non-Terrestrial
Networks}) y, en particular, de constelaciones de satélites de órbita baja (LEO,
\textit{Low Earth Orbit}), que ofrecen una pérdida de trayectoria mucho menor que los
satélites geoestacionarios y un retardo de ida y vuelta de entre 2 y 20\,ms, apropiado
para aplicaciones IoT de baja latencia~\cite{ref3, ref8}. Sin embargo, el canal LEO
introduce efectos que no existen en enlaces terrestres ---desplazamiento Doppler, retardo
variable y ventanas de visibilidad cortas--- que afectan directamente el procedimiento de
acceso aleatorio y pueden empeorar el problema de colisiones~\cite{ref8, ref7}.

Para atacar el cuello de botella del acceso, en la literatura reciente se han propuesto
esquemas de acceso aleatorio no ortogonal (NORA, \textit{Non-Orthogonal Random Access}) que,
en lugar de descartar a los dispositivos colisionados, intentan recuperarlos aprovechando
las diferencias en potencia recibida o en tiempo de llegada de sus
señales~\cite{ref5, ref6}. Variantes más avanzadas, como SARA
(\textit{Sparse Code Multiple Access-Applied Random Access}), combinan NORA con el acceso
múltiple por código escaso (SCMA, \textit{Sparse Code Multiple Access}) para separar
dispositivos usando \textit{codebooks} escasos y un detector basado en el algoritmo de paso
de mensajes (MPA, \textit{Message Passing Algorithm})~\cite{ref9, ref10}.
Estos esquemas han logrado mejoras de más del 30\,\% en capacidad en entornos terrestres,
pero su comportamiento bajo un canal satelital LEO todavía no se ha estudiado a profundidad
y no se puede deducir directamente de los resultados publicados para escenarios terrestres,
porque los efectos del canal y las ráfagas de acceso que se producen durante la ventana de
visibilidad del satélite pueden cambiar de forma importante el desempeño del esquema.

A lo anterior se suma un tercer aspecto crítico: la detección de colisiones, que
condiciona el desempeño de cualquier esquema no ortogonal. Los trabajos recientes sobre
detección basada en aprendizaje automático~\cite{ref11}, muestran que la precisión del
detector influye directamente sobre la capacidad efectiva del sistema, pero pocos estudios
han evaluado cómo se comporta esta relación cuando el canal es el de un enlace LEO.
Entender esta dependencia permite identificar qué tan robusto puede ser un esquema no
ortogonal de este tipo frente a errores de detección y bajo qué condiciones sigue valiendo
la pena usarlo.

A partir de este panorama se delimita el problema que aborda este trabajo: no está claro
cuánta capacidad del canal de acceso aleatorio se puede alcanzar al aplicar esquemas no
ortogonales como SARA sobre una red NB-IoT que usa satélites LEO como infraestructura de
acceso, ni cómo se degrada esa capacidad cuando aparecen efectos propios del enlace
satelital o errores en la detección de colisiones. Caracterizar este comportamiento aporta
información útil tanto para el diseño de futuras redes IoT no terrestres como para
orientar trabajos posteriores sobre mecanismos de acceso en NTN, un tema que la Release 18
del 3GPP identifica como un problema abierto~\cite{ref7}.
